\documentclass[11pt]{article}
\usepackage[margin=0.9in]{geometry}
\usepackage{amsmath,amssymb,amsthm,booktabs,mathrsfs}
\usepackage[T1]{fontenc}
\usepackage{lmodern,microtype}
\usepackage[colorlinks=true,linkcolor=blue,urlcolor=blue]{hyperref}
\newtheorem{theorem}{Theorem}[section]
\newtheorem{lemma}[theorem]{Lemma}
\newtheorem{proposition}[theorem]{Proposition}
\newtheorem{corollary}[theorem]{Corollary}
\theoremstyle{remark}
\newtheorem{remark}[theorem]{Remark}
\title{Toward a four-dimensional Yang--Mills construction\\
\large Notebook 34: ghost asymptotics and the physical determinant}
\author{Research working document}
\date{17 October 2026}
\begin{document}
\maketitle
\begin{abstract}
We determine the complete based-ghost contribution to the
actual leading determinant's Maxwell coefficient. An exact
returning-walk area calculation gives the unpinned lattice
heat coefficient. Counting the root residues excluded by
the pins controls the difference through a specified
heat-time interval. Through the block diffusion time
\(K^2\), the actual contribution has a fixed finite part
and an \(O(K^{-2})\) error.

The flat pinning gap is of order \(K^{-4}\), rather than
\(K^{-2}\). A pointwise Neumann Green-function estimate
proves the matching lower bound. A fourth area moment
and a killed-return heat bound then control the tail
past \(K^4\). Together they give the entire actual ghost
coefficient \(C_H\log K/(12\pi^2)+O(1)\), uniformly in
the preceding classical depth and the permitted infinite
structural systems.

This is twice the ghost logarithm of the massive reference
used in Notebook 33. The extra logarithm comes from the
interval \(K^2\) to \(K^4\) and must be compensated by the
actual vector determinant. It supplies a quantitative
test of the remaining coupled comparison, not the full
Yang--Mills existence argument.

We also prove the full actual vector short-time bound,
including the composite ordinary constraint. Its flat
derivative has a uniform counting-\(\ell^2\) bound,
and its affine magnetic jets are controlled through
the fixed-ratio hierarchy. A localized trace identity
uses the fixed coarse fiber dimension to avoid a
factor equal to the growing fine block volume.
Together with the ghost estimate, this closes the
actual ultraviolet heat interval.

Finally an exact finite determinant identity separates
the full fine and coarse gauge-orbit density factors
from a residual precision on the physical quotient.
That physical precision has a uniform flat gap of
order \(K^{-2}\). An explicit contour-gauge test shows
why the same gap cannot simply be assigned to the
unprojected contour completion. A paired scalar
reference has an exact infrared cancellation; applying
it to the actual finite identity still requires a
joint infrared limit and uniform curvature-source
estimates for the physical operator.
\end{abstract}

\section*{Inputs, conventions and scope}

The actual construction, fixed-\(K\) local infinite limits
and determinant-density identities are those of Notebooks
32 and 33. The input identities are:
\begin{enumerate}
\item On the normalized output probe
\(w_\mu(n)=t n_\nu H\), \(\mu\ne\nu\), \(\tau(H^2)=1\),
the first two actual stationary jets agree with current
affine links \(U_\mu(x)=\exp(i t x_\nu H/K^2)\).
This holds for every preceding homogeneous classical
action \(S_M\). Only these jets are used.
\item The based generator \(D_b\) has scalar domain zero
at every \(K\)-center, and \(G(t)=D_b(t)^TD_b(t)\).
The Maxwell contribution of its normalized log determinant
is \(2[t^2]\operatorname{Tr}_{0,S}\log G(t)\), where the
trace includes noncenter sites in one centered block
and their adjoint color coordinates.
\item The actual total coefficient splits into this ghost
term, the vector term
\(-[t^2]\operatorname{Tr}_{0,E}\log V_a(t)\), and a
uniformly bounded coarse term. Here
\(V_a=s_M''+D_bD_b^T+aC_K^TC_K\) on the required jets;
the constraint multiplier curvature vanishes through
the same two orders. The fixed \(a\) is the large
positive completion from Notebook 32.
\end{enumerate}

The charge factor is \(C_H=\sum_q q^2\), where \(q\)
runs over the eigenvalues of \(\operatorname{ad}_H\)
on the complexified adjoint space, with multiplicities.
There is no doubling of the real adjoint trace.
For \(SU(N)\) with \(\tau=\operatorname{Tr}_{\rm fund}/N\),
\(C_H=2N^2\). We write \(\alpha_g=C_H/(48\pi^2)\).
All estimates below are in four dimensions.

The proofs have been checked internally against these
inputs. This notebook does not assert a full
endpoint-uniform vector comparison, positive total
drift, a continuum measure, or reconstruction.



% No preamble. Prefix pgp-. Exact actual pinned ghost through fine heat time K^2.
\section{The actual ghost coefficient up to block diffusion time}
\label{pgp-section}

Work in a joint infinite-lattice structural system at fixed
\(K\), with the based centers \(K\mathbb Z^4\). Closed walks
below return in \(\mathbb Z^4\); no finite-torus winding
or continuously varying periodic flux is used.
The short-time closed-walk argument can be sharpened to
compute the whole based-ghost contribution through heat
time \(K^2\) in current fine units. This is the complete
short and intermediate interval for this term. It does
not address the remaining vector or longer-time integral.
Write
\[
 h_K^g(s)=-2[t^2]\operatorname{Tr}_{0,S}e^{-sG(t)},\qquad
 \alpha_g=\frac{C_H}{48\pi^2}.
\]
The stationary-jet identification of Notebook 33 makes
this the same first two background jets of the actual
based scalar operator for every
preceding classical depth \(M\). The field on a current
link has curvature \(tH/K^2\).

\subsection{The unpinned diagonal coefficient}
\label{pgp-unpinned-section}

Let \(X_1,\ldots,X_4\) be independent continuous-time
nearest-neighbor walks, each with jumps \(+1\) and
\(-1\) at rate one. Their joint generator is the
negative of the flat scalar lattice Laplacian.
For the selected coordinate plane, set
\[
 {\cal A}(s)=\int_0^s X_\nu(u-)\,dX_\mu(u).
\]
On a returning walk this is its oriented algebraic
area, up to the immaterial orientation sign. Define
\begin{equation}
 h_\infty^g(s)
   =C_H\,\mathbb E\left[{\cal A}(s)^2
                       {\bf1}_{\{X(s)=0\}}\right].
 \label{pgp-unpinned-definition}
\end{equation}
The \(K^4\) sites in one block cancel the curvature
factor \(K^{-4}\), so this is exactly the unpinned
ghost contribution in the final-coordinate normalization.

For \(s>0\) define
\[
 I_0(2s)=\sum_{m\ge0}\frac{s^{2m}}{(m!)^2},\qquad
 I_1(2s)=\sum_{m\ge0}\frac{s^{2m+1}}{m!(m+1)!}.
\]

\begin{proposition}[Exact lattice ghost heat coefficient]
\label{pgp-exact-heat}
For every \(s>0\),
\begin{equation}
 h_\infty^g(s)=\frac{C_H}{3}s^2e^{-8s}
                                  I_0(2s)^2I_1(2s)^2.
 \label{pgp-bessel}
\end{equation}
\end{proposition}

\begin{proof}
A single coordinate returns precisely when its two
Poisson counts of rate-one jumps agree. Its return
probability is \(p_s=e^{-2s}I_0(2s)\). Conditional
on returning, let \(m\) be the count in either
direction. The series above gives
\[
 \mathbb E[m\mid X_1(s)=0]
          =s\,\frac{I_1(2s)}{I_0(2s)}.
\]
Conditional on \(m\), the \(m\) positive jump times
and the \(m\) negative jump times are independent
uniform points in \([0,s]\). It follows by integrating
their signed point measures that, conditional on
returning,
\begin{align*}
 \mathbb E[X_1(u)X_1(v)]
   &=v_s\{\min(u,v)-uv/s\},\\
 \mathbb E[dX_1(u)dX_1(v)]
   &=v_s\{\delta(u-v)\,du-d u\,d v/s\},\qquad
 v_s=2\,\frac{I_1(2s)}{I_0(2s)}.
\end{align*}
The second equality is an identity of second-moment
measures, not a pointwise white-noise representation.
For example, at fixed \(m\) its diagonal mass is
\(2m\,du/s\); the off-diagonal density is
\(-2m\,du\,dv/s^2\).

Conditioning the four coordinates to return preserves
their mutual independence. There are almost surely
no simultaneous jumps in different coordinates.
Using the second-moment measure of \(X_\mu\) and
the bridge covariance of \(X_\nu\), therefore,
\begin{align*}
 \mathbb E[{\cal A}(s)^2\mid X(s)=0]
 &=v_s^2\left\{
      \int_0^s(u-u^2/s)\,du
       -\frac1s\int_0^s\int_0^s
                         \{\min(u,v)-uv/s\}\,du\,dv\right\}\\
 &=\frac{v_s^2s^2}{12}
   =\frac{s^2}{3}\left(\frac{I_1(2s)}{I_0(2s)}\right)^2.
\end{align*}
Multiplication by \(p_s^4\) and \(C_H\) proves
\eqref{pgp-bessel}. All counts have finite moments,
so these expectation interchanges are justified.
\end{proof}

\begin{lemma}[Bounds with an integrable large-time error]
\label{pgp-asymptotics}
There is a constant \(B\), depending only on the color
normalization, such that
\begin{align}
 &0\le h_\infty^g(s)\le B\min(s^4,1),\qquad s>0,
                                      \label{pgp-uniform-bound}\\
 &|h_\infty^g(s)-\alpha_g|\le B/s,\quad
   |s(h_\infty^g)'(s)|\le B,\qquad s\ge1.
                                      \label{pgp-asymptotic-bound}
\end{align}
\end{lemma}

\begin{proof}
The series give the small-\(s\) estimate. For \(s\ge1\)
put \(a_s=e^{-2s}I_0(2s)\), \(b_s=e^{-2s}I_1(2s)\).
Fourier coefficient extraction from
\(e^{s(e^{i\theta}+e^{-i\theta})}\) gives
\[
 a_s=\frac1{2\pi}\int_{-\pi}^{\pi}
                   e^{-2s(1-\cos\theta)}\,d\theta,\qquad
 b_s=\frac1{2\pi}\int_{-\pi}^{\pi}
                   e^{-2s(1-\cos\theta)}\cos\theta\,d\theta.
\]
On this interval, with \(\phi(\theta)=2(1-\cos\theta)\),
\[
 \frac{4}{\pi^2}\theta^2\le\phi(\theta)\le\theta^2,\qquad
 |\phi(\theta)-\theta^2|\le\theta^4/12.
\]
Comparison with the full Gaussian integral gives
\[
 a_s=(4\pi s)^{-1/2}+O(s^{-3/2}).
\]
Indeed the integrand error is bounded by
\(Cs\theta^4e^{-c s\theta^2}\), whose integral is
\(O(s^{-3/2})\), and the Gaussian tail is smaller.
The factor \(\cos\theta-1=O(\theta^2)\) gives the
same asymptotic formula for \(b_s\).
The integral formulas also give
\[
 |a_s|+|b_s|\le Cs^{-1/2},\qquad
 |a_s'|+|b_s'|\le Cs^{-3/2}.
\]
Substitute these estimates into
\(h_\infty^g=(C_H/3)s^2a_s^2b_s^2\).
Its limiting value is \(C_H/(48\pi^2)\); the
relative product error is \(O(s^{-1})\).
Differentiating the product with the last bounds
proves the stated derivative bound.
\end{proof}

\subsection{Pinning changes this prefix by a vanishing amount}
\label{pgp-pinning-section}

The centers form \(K\mathbb Z^4\). A centered block
contains exactly one representative of each residue
class modulo \(K\) in each coordinate. For a closed
walk shape \(\gamma\) of length \(n\ge1\) rooted at zero,
let \(r_K(\gamma)\) be the number of these root residues
whose translated walk meets a center. Each vertex
of the cyclic walk excludes one root residue at most,
so
\begin{equation}
               0\le r_K(\gamma)\le\min(n,K^4).
 \label{pgp-root-count}
\end{equation}
This includes roots which themselves are centers.
Algebraic area is invariant under translating a closed
walk.

\begin{proposition}[Uniform pinning comparison through time \(K^2\)]
\label{pgp-pinning}
For every \(s>0\),
\begin{equation}
 0\le h_\infty^g(s)-h_K^g(s)
 \le K^{-4}\{s(h_\infty^g)'(s)+8s h_\infty^g(s)\}.
 \label{pgp-pointwise-pinning}
\end{equation}
In particular, for \(K\ge1\),
\begin{equation}
 \int_0^{K^2}|h_K^g(s)-h_\infty^g(s)|\,\frac{ds}{s}
                                      \le B'K^{-2}
 \label{pgp-prefix-error}
\end{equation}
with \(B'\) independent of \(K,M\).
\end{proposition}

\begin{proof}
The full closed-walk expansion and the actual
center-avoiding expansion have nonnegative squared-area
coefficients. Their difference is exactly
\[
 \frac{C_H}{K^4}e^{-8s}
   \sum_{n\ge1}\frac{s^n}{n!}
       \sum_{\gamma:0\to0,\ |\gamma|=n}
                         r_K(\gamma){\cal A}(\gamma)^2 .
\]
Inserting \(r_K(\gamma)\le n\) gives
\[
 \frac1{K^4}e^{-8s}
           s\frac{d}{ds}\{e^{8s}h_\infty^g(s)\},
\]
which is \eqref{pgp-pointwise-pinning}. The quantity
in braces on its right is nonnegative by this
generating-series representation; no monotonicity
of \(h_\infty^g\) alone is assumed.
Integrating from zero to \(T\) yields
\[
 \int_0^T(h_\infty^g-h_K^g)\,\frac{ds}{s}
 \le K^{-4}\left\{h_\infty^g(T)
                       +8\int_0^T h_\infty^g(s)\,ds\right\},
\]
since \(h_\infty^g(0)=0\).
The uniform bound \eqref{pgp-uniform-bound}, with
\(T=K^2\), proves \eqref{pgp-prefix-error}.
\end{proof}

\begin{theorem}[Actual ghost prefix with a fixed finite part]
\label{pgp-prefix-theorem}
Define the finite constant
\[
 \gamma_g=\int_0^1 h_\infty^g(s)\,\frac{ds}{s}
       +\int_1^\infty\{h_\infty^g(s)-\alpha_g\}\,\frac{ds}{s}.
\]
Then the actual based-ghost contribution satisfies
\begin{equation}
 \boxed{\quad
 \int_0^{K^2}h_K^g(s)\,\frac{ds}{s}
     =\frac{C_H}{24\pi^2}\log K+\gamma_g+O(K^{-2}).
 \quad}
 \label{pgp-prefix-law}
\end{equation}
The error constant is independent of the preceding
classical depth \(M\), \(K\), and the permitted joint
structural system.
\end{theorem}

\begin{proof}
Both integrals defining \(\gamma_g\) converge by
Lemma~\ref{pgp-asymptotics}. The tail of the second
integral from \(K^2\) to infinity is \(O(K^{-2})\).
Thus the same formula first holds with \(h_\infty^g\).
Proposition~\ref{pgp-pinning} then proves it for the
actual based ghost. Its independence of the effective
action follows from the previously established
affine stationary-jet identification.
\end{proof}

There is also the pointwise intermediate-scale bound
\begin{equation}
 \left|h_K^g(K^2u)-\alpha_g\right|
 \le B''\{(K^2u)^{-1}+u/K^2\},
                         \qquad K^{-2}\le u\le1.
 \label{pgp-band-bound}
\end{equation}
Indeed \eqref{pgp-pointwise-pinning} and the derivative
bound give an error at most \(Cs/K^4\) for \(s\ge1\),
and the unpinned error is \(O(s^{-1})\).
This verifies the actual ghost part of the middle
heat-time comparison with exponent one.

The remaining vector coefficient on that interval
must supply \(10\alpha_g\) in the continuum convention
of Notebook 33. The estimate above does not bound
\(\int_{K^2}^\infty h_K^g(s)\,ds/s\) uniformly.
At those longer times the pinning gap and coupled
vector contribution must be retained. In particular,
\eqref{pgp-prefix-law} is a theorem about a specified
heat-time prefix, not a formula for the entire actual
determinant or its sign.


% No preamble. Prefix pgl-. Entire pinned ghost coefficient; vector still needed.
\section{The pinning scale and the entire based-ghost logarithm}
\label{pgl-section}

The pinning scale is longer than the block diffusion
time. We now control it sharply enough to determine
the whole actual based-ghost Maxwell contribution.
All statements retain the infinite-system and formal
jet conventions of Section~\ref{pgp-section}.

\subsection{The correct flat pinned gap}
\label{pgl-gap-section}

\begin{lemma}[Four-dimensional point-pinned Poincare bound]
\label{pgl-gap}
Let \(G_K(0)\) be the flat scalar lattice Laplacian on
functions vanishing at every site of \(K\mathbb Z^4\).
There is a constant \(c>0\), independent of \(K\), such
that
\begin{equation}
                     G_K(0)\ge cK^{-4}I.
 \label{pgl-gap-bound}
\end{equation}
The statement holds on compatible finite periodic
volumes and on the infinite lattice.
\end{lemma}

\begin{proof}
It suffices to prove a block inequality using only
edges internal to each centered cube. For \(K=1\)
the domain is zero; assume \(K\ge2\).
The Neumann graph Laplacian of a path with \(K\)
vertices has eigenvalues
\[
 \lambda_j=4\sin^2\frac{\pi j}{2K},\qquad
                  0\le j<K.
\]
Its normalized cosine eigenvectors have squared
values at most \(2/K\). Products give an eigenbasis
\(\psi_{\boldsymbol j}\) in four dimensions, with
\[
 |\psi_{\boldsymbol j}(x)|^2\le16K^{-4},\qquad
 \lambda_{\boldsymbol j}
       =\sum_{\ell=1}^4\lambda_{j_\ell}
                          \ge c|\boldsymbol j|^2/K^2.
\]
Consequently the point diagonal of the zero-mean
Neumann inverse is uniformly bounded:
\[
 \sum_{\boldsymbol j\ne0}
       \frac{|\psi_{\boldsymbol j}(x)|^2}
            {\lambda_{\boldsymbol j}}
 \le CK^{-2}
       \sum_{\substack{\boldsymbol j\in\{0,\ldots,K-1\}^4\\
                        \boldsymbol j\ne0}}
                              |\boldsymbol j|^{-2}\le C.
\]
The last sum is \(O(K^2)\), by shells in four dimensions.

Write a block field as \(\phi=\bar\phi+u\), with
\(u\) of mean zero. Its internal edge energy is
denoted by \(E_B(\phi)\). Spectral Cauchy--Schwarz
gives \(|u(x)|^2\le C E_B(\phi)\), and the first
nonzero Neumann eigenvalue gives
\(\|u\|_2^2\le CK^2 E_B(\phi)\).
At the pinned center, \(\bar\phi=-u(x_{\rm center})\).
It follows that
\[
 \|\phi\|_{2,B}^2=\|u\|_{2,B}^2+K^4|\bar\phi|^2
                                  \le CK^4 E_B(\phi).
\]
Sum over disjoint blocks. Internal edges are a
subset of all edges, proving \eqref{pgl-gap-bound}.
The infinite statement follows first on finitely
supported fields and then by closure.
\end{proof}

Together with the constant-off-centers trial function
from Notebook 32, this gives the sharp order \(K^{-4}\)
of the bottom of the flat based spectrum. The finite
periodic upper test gives at most \(8/(K^4-1)\);
the same order follows in infinite volume by large
cutoffs. It does not give a mass of order \(K^{-1}\).

\subsection{A killed-return bound and an area fourth moment}
\label{pgl-tail-inputs-section}

Let \(p_s^K(x,y)\) be the heat kernel killed on visiting
a center. At the flat background it is the kernel of
\(e^{-sG_K(0)}\). For \(s\ge1\), free-kernel domination
and \eqref{pgl-gap-bound} imply
\begin{equation}
 p_s^K(x,x)\le Cs^{-2}e^{-c s/K^4},
                         \qquad x\notin K\mathbb Z^4.
 \label{pgl-killed-return}
\end{equation}
Indeed split \(e^{-sG_K/2}\delta_x\) into a time
\(s/4\) contraction followed by time \(s/4\)
smoothing:
\[
 p_s^K(x,x)
 =\|e^{-sG_K/2}\delta_x\|_2^2
 \le e^{-c s/(2K^4)}p_{s/2}^K(x,x)
 \le e^{-c s/(2K^4)}p_{s/2}^{\rm free}(0,0).
\]
The last free return is \(O(s^{-2})\) by the
one-dimensional integral estimates already proved.
Absorb the harmless factor two in \(c\).

\begin{lemma}[Fourth area moment of the returning walk]
\label{pgl-area-fourth}
For \(s\ge1\), the free four-dimensional walk satisfies
\begin{equation}
 \mathbb E[{\cal A}(s)^4\mid X(s)=0]\le Cs^4,\qquad
 \mathbb E[{\cal A}(s)^4{\bf1}_{\{X(s)=0\}}]\le Cs^2.
 \label{pgl-fourth-bound}
\end{equation}
\end{lemma}

\begin{proof}
For a returning one-dimensional coordinate, its
balanced jump count \(m\) satisfies the exact identity
\[
 \mathbb E[m^2\mid X_1(s)=0]
 =\frac{\sum_{m\ge1}m^2s^{2m}/(m!)^2}{I_0(2s)}
 =s^2.
\]
Thus \(\mathbb E m\le s\) as well. Given the bridge
path \(Y=X_\nu\) and the count \(m\) for \(X_\mu\),
its area is
\[
 {\cal A}(s)=\sum_{i=1}^m Y(U_i)-\sum_{i=1}^m Y(V_i),
\]
where the \(U_i,V_i\) are independent uniforms in
\([0,s]\), independent of \(Y\). The deterministic
mean terms cancel. The fourth moment formula for a
sum of independent centered variables and Jensen's
inequality give
\[
 \mathbb E[{\cal A}(s)^4\mid m,Y]
 \le C(m+m^2)\,\frac1s\int_0^s|Y(u)|^4\,du.
\]
For clarity, centering \(Y(U_i)\) changes its fourth
moment by at most a fixed factor; the square of its
second moment is at most its fourth moment.

Given the balanced count \(l\) of \(Y\), the value
\(Y(u)\) is the difference of two independent
\(\operatorname{Binomial}(l,u/s)\) variables.
The same centered-sum formula bounds its fourth
moment by \(C(l+l^2)\), uniformly in \(u\).
Averaging counts therefore gives
\[
 \mathbb E[|Y(u)|^4\mid Y(s)=0]\le C(s+s^2)\le Cs^2.
\]
The two coordinate bridges and their counts are
independent. Combining the last bounds proves
the conditional \(Cs^4\) estimate. The four-dimensional
return probability is \(p_s^4\le Cs^{-2}\), proving
the unconditional bound.
\end{proof}

\begin{proposition}[Uniform ghost tail past the pinning scale]
\label{pgl-tail}
For \(s\ge1\),
\begin{equation}
 0\le h_K^g(s)\le C e^{-c s/K^4}.
 \label{pgl-ghost-tail-pointwise}
\end{equation}
In particular,
\begin{equation}
           \int_{K^4}^\infty h_K^g(s)\,\frac{ds}{s}\le C,
 \label{pgl-tail-integral}
\end{equation}
uniformly in \(K,M\).
\end{proposition}

\begin{proof}
At each noncenter starting site \(x\), the squared-area
walk formula has the probabilistic form
\[
 h_K^g(s)=\frac{C_H}{K^4}
 \sum_{\substack{x\ {\rm noncenter}\\x\ {\rm in\ block}\ 0}}
 \mathbb E_x\left[{\cal A}(s)^2
    {\bf1}_{\{X(s)=x,\ X\ {\rm avoids\ centers}\}}\right].
\]
Translate the area coordinate to \(x\), which changes
no area on a closed walk. Cauchy--Schwarz bounds each
expectation by
\[
 \left(\mathbb E[{\cal A}(s)^4
                       {\bf1}_{\{X(s)=0\}}]\right)^{1/2}
                  \left(p_s^K(x,x)\right)^{1/2}
 \le C e^{-c s/K^4},
\]
using \eqref{pgl-killed-return} and
\eqref{pgl-fourth-bound}. The block has at most \(K^4\)
sites, canceling the prefactor. Integrate with
\(s=K^4v\), \(v\ge1\), to obtain
\eqref{pgl-tail-integral}.
\end{proof}

\subsection{The additional logarithm that the vector must cancel}
\label{pgl-full-section}

\begin{theorem}[Entire actual based-ghost coefficient]
\label{pgl-full}
There is a constant \(C\), independent of \(M,K\) and
the permitted joint structural system, such that
\begin{equation}
 \left|\int_0^\infty h_K^g(s)\,\frac{ds}{s}
                      -4\alpha_g\log K\right|\le C.
 \label{pgl-full-law}
\end{equation}
More specifically,
\begin{equation}
 \int_{K^2}^{K^4} h_K^g(s)\,\frac{ds}{s}
                         =2\alpha_g\log K+O(1),
 \label{pgl-extra-log}
\end{equation}
and the tail past \(K^4\) is uniformly bounded.
\end{theorem}

\begin{proof}
The proof of Proposition~\ref{pgp-pinning} gave, for
every \(T>0\),
\[
 0\le\int_0^T(h_\infty^g-h_K^g)\,\frac{ds}{s}
 \le K^{-4}\left\{h_\infty^g(T)
                    +8\int_0^T h_\infty^g(s)\,ds\right\}.
\]
Set \(T=K^4\). The right side is uniformly bounded,
while the unpinned integral is
\(\alpha_g\log(K^4)+\gamma_g+O(K^{-4})\).
This proves the full prefix law with an \(O(1)\)
error. Add Proposition~\ref{pgl-tail} to get
\eqref{pgl-full-law}. Subtracting
\eqref{pgp-prefix-law} gives \eqref{pgl-extra-log}.
\end{proof}

Recall explicitly that
\[
 \int_0^\infty h_K^g(s)\,\frac{ds}{s}
   =2[t^2]\operatorname{Tr}_{0,S}\log G(t)
   =\beta\bigl(\mathcal N_0\log\det G\bigr),
\]
with the determinant-density and Maxwell normalization
of Notebook 33. Thus the displayed factor two is the
affine coefficient extraction, not an extra determinant.
Its coefficient \(4\alpha_g=C_H/(12\pi^2)\) is
twice the \(2\alpha_g=C_H/(24\pi^2)\) of the
continuum ghost with mass \(K^{-1}\) in current
fine units and the same ultraviolet scale.
This difference is an actual pinning effect in the
separated determinant; it is not a failure of the
normalization.

In particular, a uniform estimate on the separated
ghost tail starting at \(K^2\) is false. Its positivity
and \eqref{pgl-extra-log} prove a logarithmic lower
bound on that tail. The total reference coefficient
is \(22\alpha_g\log K\); hence a bounded comparison
for the full actual determinant requires the vector
contribution to be \(18\alpha_g\log K+O(1)\), since
the coarse determinant is bounded. Equivalently,
the vector must compensate the additional ghost
logarithm relative to its massive reference.
No such vector estimate is established by this
scalar calculation.


% No preamble. Prefix vst-. Actual composite-constraint vector UV bound.
\section{A uniform short-time estimate for the actual vector operator}
\label{vst-section}

Fix the primitive odd ratio \(\ell\ge3\), the actual selected analytic
source mean, and \(a>0\). Let \(K=\ell^k\), and let \(M\) be any
preceding homogeneous depth. We work in an actual joint infinite
structural system, with current-link angle counting coordinates.
All affine substitutions below are polynomial jets at zero.
No continuously varying periodic magnetic flux is used.

Choose \(\mu\ne\nu\) and a Cartan element \(H\) with
\(\tau(H^2)=1\), and write
\[
 U_\mu(t;x)=\exp(i t x_\nu H/K^2),\qquad U_\lambda(t;x)=1
                       \quad(\lambda\ne\mu).
\]
The stationary-jet identification established in Notebook 33 gives,
through order two, the actual vector family
\begin{equation}
 \begin{split}
 B_M(t)&=s_M''(U(t))+D_b(t)D_b(t)^T,\\
 V_{M,K}(t)&=B_M(t)+aC_K(t)^TC_K(t).
 \end{split}
 \label{vst-operators}
\end{equation}
Here \(s_M''\) is the actual left-relative Hessian,
\(D_b\) has scalar values zero at the \(K\)-centers, and \(C_K\)
is the left-relative derivative of the composite ordinary map
\(M_\ell^k\). In particular, \(s_M''\) is not replaced by the
bare Wilson Hessian when \(M>1\).

The inputs used from the classical construction are uniform
current-coordinate exponential bounds for \(q_M=s_M''(1)\),
the complete scalar-action support/Cauchy bounds through order
four on a fixed complex product ball, and the exact local gauge
covariance of the relative Hessian. The local coordinate and
color fibers have fixed dimension. The actual primitive
averaging formula supplies finite-range analytic derivative
matrices. Its composite output row reads only the two final
endpoint blocks. These inputs are uniform in preceding depth
and compatible regulator volume.

\begin{theorem}[Actual vector ultraviolet bound]
\label{vst-theorem}
There is a constant \(C\), depending on the fixed primitive
data, \(a\), the classical analytic bounds and the color group,
but not on \(M,K\) or the selected structural system, such that
\begin{equation}
 \int_0^1\left|[t^2]\operatorname{Tr}_{0,E}
                         e^{-sV_{M,K}(t)}\right|\frac{ds}{s}
 \le C.
 \label{vst-main}
\end{equation}
The trace is over fine bonds with initial site in block zero,
including color. Together with the actual ghost ultraviolet
bound, this proves the full short-time estimate on
\(0<s<1\), or \(0<u<K^{-2}\) after \(s=K^2u\).
\end{theorem}

\subsection{Uniform flat bounds for the full composite constraint}
\label{vst-flat-section}

Let \(D_r\) be the scalar forward difference on the lattice
of spacing \(r\), without division by \(r\). Define the
straight angle average
\begin{equation}
 (A_r v)(y,\lambda)=r^{-d}
       \sum_{x\in B_r(y)}\sum_{h=0}^{r-1}
                              v_\lambda(x+h e_\lambda).
 \label{vst-straight}
\end{equation}
Let \(b_\ell\) be the scalar mean of internal primitive
root-contour sums. The actual arithmetic differential is
\begin{equation}
 C_\ell=A_\ell-D_\ell b_\ell,\qquad
 C_\ell D_1=D_\ell R_\ell,
 \label{vst-primitive}
\end{equation}
where \(R_\ell\) restricts scalar values to the new centers.
Tiling centered blocks and concatenating strings gives
\(A_\ell^{(r)}A_r=A_{\ell r}\). Consequently
\begin{equation}
 \begin{split}
 C_K&=A_K-D_Kb_K,\\
 b_{\ell r}&=b_\ell^{(r)}A_r+R_\ell b_r,\\
 b_K&=\sum_{j=0}^{k-1}
       R_{\ell^{j+1}\to K}\,b_\ell^{(\ell^j)}A_{\ell^j}.
 \end{split}
 \label{vst-recursion}
\end{equation}
Each scalar restriction in the last line is a counting
\(L^2\) contraction. At a final center it merely retains
the same nested center; it does not sum all intermediate
centers.

The absolute row sums of \(A_r\) equal \(r\), and its
absolute column sums equal \(r^{1-d}\). For a column,
the two neighboring longitudinal output intervals have
complementary multiplicities summing to \(r\), and the
transverse output block is unique. Thus
\begin{equation}
                \|A_r\|_{2\to2}\le r^{1-d/2}.
 \label{vst-straight-norm}
\end{equation}
These counting identities retain their multiplicities
on compatible periodic graphs. Since \(b_\ell\) is a
fixed finite-range operator and
\(\|D_K\|_{2\to2}\le2\sqrt d\), in \(d=4\)
\begin{equation}
 \|C_K\|_{2\to2}
 \le K^{-1}+
       \frac{2\sqrt d\,\|b_\ell\|_{2\to2}}{1-\ell^{-1}}.
 \label{vst-full-flat}
\end{equation}
Every consecutive flat composite segment therefore has
norm at most one fixed \(C_0\ge1\), including empty segments.
This is an \(O(1)\) bound for the full ordinary map, not
the generally false \(O(K^{-1})\) assertion.

For localized information, a row of \(A_r\) has at most
\(2r^d\) entries, each at most \(r^{1-d}\), so
\[
              \|A_r(c,\cdot)\|_2\le\sqrt2\,r^{-1}
                         \quad(d=4).
 \]
Equation \eqref{vst-recursion} gives uniformly bounded
rows of \(b_K\), and hence of \(C_K\).
The two-block dependency also shows that each input
bond occurs in only a dimension-dependent number of
different final output rows. This gives a second direct
row-and-incidence proof of the full bound.

\subsection{Actual affine jets of the nonlinear composite derivative}
\label{vst-magnetic-section}

Fix first a final row with endpoint blocks near zero.
At stage \(j\), put \(r=\ell^j\). The exact affine
averaging identity gives the intermediate background
angle, in that stage's lattice coordinates,
\[
                     w_j(z)=t(r^2/K^2)z_\nu H.
\]
The exact nested dependency geometry confines every
needed coordinate to \(|z|\le c_\ell K/r\). Therefore
\begin{equation}
                  \|w_j\|_{\rm needed}
                            \le c_\ell |t|r/K.
 \label{vst-stage-smallness}
\end{equation}
At each stage independently, extend the needed values
to a globally bounded field with this same bound, by
clipping the coordinate outside a larger dependency box.
The extensions need not be successive blockings.
Every derivative entry contributing to the selected
product row agrees with the actual affine stage, which
is the only use of these extensions.

Let \(F_j\) be the flat primitive derivative and
\(E_j(t)\) its change at the extended background.
The fixed primitive analytic and incidence bounds
give, on one disk \(|t|<\rho_0\),
\begin{equation}
                  \|E_j(t)\|_{2\to2}
                          \le C_\ell |t|\ell^j/K.
 \label{vst-stage-error}
\end{equation}
Expand the derivative product by the subset of stages
where \(E_j\) is used. A term with \(n\) marked stages
has \(n+1\) intervening flat segments, each bounded by
\(C_0\). Hence
\begin{align}
 \left\|\prod_j(F_j+E_j(t))\right\|
 &\le C_0\prod_j(1+C_0\|E_j(t)\|)\notag\\
 &\le C_0\exp\!\left(C C_0|t|
                         \sum_{j<k}\ell^j/K\right)\notag\\
 &\le C_0\exp(C C_0|t|/(\ell-1)).
 \label{vst-product-bound}
\end{align}
The actual selected row consequently has a uniform
analytic \(L^2\) bound. Cauchy estimates on a smaller
fixed disk bound its Taylor coefficients through order two.

For a row at coarse position \(y\), translating the
affine field adds a constant Cartan potential. Remove
it by the site gauge \(\exp(i t y_\nu x_\mu H/K)\).
Exact gauge covariance expresses the row using its
centered counterpart and input/output adjoint frames.
On its two \(K\)-blocks the frame derivatives through
order two are bounded by a fixed polynomial in \(|y|\),
independently of \(K\). Writing
\(C_K(t)=\sum_j t^j C^{[j]}\), one may take
\begin{equation}
        \|C^{[j]}(y,\cdot)\|_2
                  \le C_j(1+|y|)^{2j},
                  \qquad j=0,1,2.
 \label{vst-row-jets}
\end{equation}
Their support remains in the same two endpoint blocks.
This is a statement about polynomial jets; it does not
assert a bounded analytic operator at a nonzero unbounded
affine field.

Group current bonds by their initial-site \(K\)-block,
and final bonds by their initial coarse site. The latter
fiber has fixed dimension \(r_c=d\dim\mathfrak g\).
Equation \eqref{vst-row-jets} and finite incidence give
finite-range block operator bounds with polynomial
growth for \(C^{[j]}\) and their transposes.
In particular \(C^{[0]}\) has a uniform weighted block
Schur norm for every fixed block exponent.

\subsection{The base vector before the constraint insertion}
\label{vst-base-section}

At zero,
\[
                 B_0=q_M+D_bD_b^T
\]
is positive and has a uniform exponentially weighted
fine-lattice Schur norm. A spectral gap is unnecessary
on \(0\le s\le1\): the Schur-algebra power series bounds
its semigroup uniformly with exponential matrix decay.

The actual scalar-action bounds through order four
give uniform exponentially rooted kernels for the
first two background derivatives of \(B\).
The nearest-neighbor based-gradient term has the same
property uniformly in the center spacing.
For each fine diagonal bond rooted at \(x\), use a
Cartan gauge to subtract the constant potential
\(t x_\nu H/K^2\). The diagonal color heat trace is
unchanged. Its source field is now
\(t(z_\nu-x_\nu)H/K^2\). Each linear magnetic vertex
has a factor \(K^{-2}\), and each quadratic vertex
has \(K^{-4}\); uniform exponential kernels absorb
the remaining coordinate polynomials.

The one quadratic insertion and the ordered two linear
insertions in the heat expansion therefore give
\begin{equation}
 \left|[t^2]\operatorname{tr}_{\mathfrak g}
                         (e^{-sB(t)})_{xx}\right|
        \le CK^{-4}(s+s^2)e^{Cs},\qquad 0<s\le1.
 \label{vst-base-diagonal}
\end{equation}
The local gauge identities extend to these polynomial
jets by their exponentially summable kernels.
There are \(dK^4\) bond roots per block, so
\begin{equation}
 \left|[t^2]\operatorname{Tr}_{0,E}e^{-sB(t)}\right|
                  \le C(s+s^2)e^{Cs}.
 \label{vst-base-trace}
\end{equation}
Its \(ds/s\) integral is uniformly finite.

\subsection{Uniform block estimates without tracing the large fiber}
\label{vst-block-section}

The coefficients \(B^{[j]}\), \(j\le2\), also have
uniform exponentially decaying \(K\)-block operator
kernels, with fixed polynomial growth in the block
indices for \(j>0\). Fine positions in block \(X\)
have magnitude at most \(CK(1+|X|)\), while the
magnetic source coefficients contain \(K^{-2j}\).
Fine Schur bounds control the operator norm of each
block without its dimension as a prefactor.
For distant blocks the fine distance is at least
\(K(d_{\rm block}(X,Y)-d)\); a smaller fixed block
exponent is therefore summable uniformly in \(K\).
The finitely many nearby block offsets are controlled
by the full fine operator norm.

Set \(A(t)=aC_K(t)^TC_K(t)\). The preceding finite-range
block bounds apply to \(A^{[j]}\) and to
\(V^{[j]}=B^{[j]}+A^{[j]}\).
For \(E=B,V\), the block power series at zero and the
one/two insertion formulas imply, for \(j\le2\),
\begin{equation}
 \big\|([t^j]e^{-sE(t)})_{XY}\big\|
 \le C(1+|X|+|Y|)^p e^{-\kappa|X-Y|},
                         \qquad 0\le s\le1.
 \label{vst-semigroup-block}
\end{equation}
All constants are uniform in \(M,K\).
For positive \(j\) the usual factors \(s\) and
\(s+s^2\) can also be retained. Products of polynomial
kernels are controlled by spending a smaller block
exponent. No trace over a \(K^4\)-dimensional fiber
has yet been taken.

\subsection{The constraint correction has fixed coarse trace rank}
\label{vst-rank-section}

Let \(P_0\) project onto current bonds rooted in block zero.
The exact Duhamel difference is
\begin{equation}
 e^{-sV(t)}-e^{-sB(t)}
   =-a\int_0^s e^{-(s-u)V(t)}
              C_K(t)^TC_K(t)e^{-uB(t)}\,du.
 \label{vst-difference}
\end{equation}
At each coefficient through degree two, cyclicity gives
\begin{equation}
 \operatorname{Tr}_{\rm fine}
           (P_0E_VC^TCE_B)
       =\operatorname{Tr}_{\rm coarse}
           (CE_BP_0E_VC^T).
 \label{vst-local-cyclicity}
\end{equation}
The projector \(P_0\) remains in this identity. The
right side sums all coarse diagonal blocks; it is
not identified with an arbitrary block-zero coarse
trace density.

For each coarse root \(Y\), its diagonal block on the
right is a product through the fine block-zero Hilbert
space. The outer fiber has dimension \(r_c\), independent
of \(K\). Equations \eqref{vst-row-jets} and
\eqref{vst-semigroup-block} bound the coefficients through
degree two of the two maps by
\[
                     C(1+|Y|)^p e^{-\kappa'|Y|}.
\]
The absolute trace of their product is at most \(r_c\)
times the product of these norms. Summing \(Y\) is
uniformly finite. Expanding the quadratic coefficient
in \eqref{vst-local-cyclicity}, and then integrating
the Duhamel time, proves
\begin{equation}
 \left|[t^2]\operatorname{Tr}_{0,E}
          (e^{-sV(t)}-e^{-sB(t)})\right|
                    \le Cs,\qquad 0<s\le1.
 \label{vst-rank-bound}
\end{equation}
One can first use finite boxes and compact sources.
The displayed rooted sums are absolute and justify
the limit. For fixed \(K\), the finite \(P_0\) also
makes ordinary trace cyclicity legitimate; the
estimate itself uses the small outer coarse fiber
to avoid a \(K^4\) factor. Complex coefficients use
the holomorphic bilinear transpose, which agrees
with the adjoint at the real flat point.

Combining \eqref{vst-rank-bound} with
\eqref{vst-base-trace} and integrating \(ds/s\)
proves Theorem~\ref{vst-theorem}.

\subsection{What this estimate does and does not settle}
\label{vst-scope}

The argument uses the actual composite fixed-\(\ell\)
hierarchy. A direct-\(K\) contour average can have
different raw row bounds and was not substituted.
The full ordinary map needs only the proved
\(O(1)\) counting bound; sampled-center gauge
directions forbid assuming \(O(K^{-1})\) for it.

The final affine statement is in the infinite
structural system. The uniform finite-regulator
inputs and local limits do not turn it into a
continuous magnetic-flux formula on an arbitrary
torus. The semigroup estimates here use \(e^{Cs}\)
on bounded time, not a uniform large-time gap.
They prove neither the continuum coefficient on
intermediate times nor the cancellation of the
soft based-ghost contribution at large times.
A combined signed-integral comparison remains
an independent problem.


\section{An exact finite physical quotient for the determinant}
\label{pqd-section}

The identities in this section concern the actual finite homogeneous
small-field germ. They remove the periodically pinned gauge variables
before estimation. They do not identify the separate massless factors
below with infinite-volume affine magnetic traces.

\subsection{The actual spaces and the residual Hessian}
\label{pqd-spaces}

Fix an old depth $M$, a composite endpoint $K=\ell^k$, and a compatible
finite current torus. All matrices initially use real, trace-orthonormal
color coordinates, current bond angles, and counting measure. Let $E$
and $E_c$ be the fine/current and final coarse bond spaces. Fix one final
coarse site $o$ as a common global gauge base. Let $S_f,S_c$ denote the
fine and coarse scalar site spaces whose value at $o$ is zero. The
noncenter scalar space $S_b$ vanishes at every final coarse center.
There is a coordinate splitting
\begin{equation}
 S_f=S_b\oplus S_c,\qquad r:S_f\longrightarrow S_c,
 \label{pqd-scalar-split}
\end{equation}
where $r$ is restriction to coarse centers; the coordinate extension of
$S_c$ is zero at noncenters.

At the actual fine saddle $U$ and target $W$, write $D_f,D_c$ for the
left-relative gauge generators with these bases, and $D_b$ for the
restriction of $D_f$ to $S_b$. Each generator with a global base is
injective, by propagating a zero scalar value along lattice paths.
The actual constraint differential $C:E\to E_c$ is onto, by its
invertible pivot block, and exact gauge covariance gives
\begin{equation}
 CD_f=D_c r,\qquad CD_b=0.\label{pqd-covariance}
\end{equation}
Use the full relative-coordinate Lagrangian and coarse Hessians
\begin{equation}
 L=s''-\lambda\cdot f'',\qquad q=s_{\rm eff}'',\qquad
 \mathcal H=L-C^TqC.\label{pqd-residual}
\end{equation}
The actual minimum-value comparison and its envelope identities give
\begin{equation}
 \mathcal H\ge0,\qquad \mathcal H D_f=0,\qquad
 CH=I,\qquad \mathcal H H=0,\label{pqd-residual-identities}
\end{equation}
where $H$ is the derivative of the actual minimum section.
The constraint-current term $\lambda\cdot f''$ is retained throughout.
No on-shell hypothesis is required.

Let $T$ be the actual contour-logarithm differential. Its normalization
is $TD_b=-I$. Positivity of the actual retained precision implies that
$\mathcal H$ on $\ker C$ has nullspace exactly $\operatorname{ran}D_b$
and is positive on its quotient. Indeed $\mathcal H=L$ on $\ker C$,
and the contour component can be removed uniquely along $D_b$.

Define
\begin{equation}
 G_f=D_f^TD_f,\quad G_c=D_c^TD_c,\quad
 P_f=I-D_fG_f^{-1}D_f^T,\quad
 P_c=I-D_cG_c^{-1}D_c^T.\label{pqd-projections}
\end{equation}
Choose orthonormal column matrices $U_f,U_c$ spanning the ranges of
$P_f,P_c$, and put, for a fixed $a>0$,
\begin{equation}
 A=U_c^TCU_f,\qquad \mathcal H_\perp=U_f^T\mathcal H U_f,
 \qquad Q_\perp=\mathcal H_\perp+aA^TA.
 \label{pqd-physical-precision}
\end{equation}
These determinants do not depend on the orthonormal bases.
The map $A$ is onto: $P_fHU_c$ is a right inverse after passing to
physical coordinates, since $P_cCD_f=0$. Its range is annihilated by
$\mathcal H$. It follows that $Q_\perp>0$ and
\begin{equation}
 A Q_\perp^{-1}A^T=a^{-1}I.\label{pqd-compression}
\end{equation}
For example, if $H_\perp=U_f^TP_fHU_c$, then
$AH_\perp=I$, $\mathcal H_\perp H_\perp=0$, and
$Q_\perp H_\perp=aA^T$, proving the last identity.

\subsection{The exact determinant and orbit-density factors}
\label{pqd-identity-section}

Write $\mathcal N_0F=F(W)-F(1)$, with the determinant branches continued
from the positive flat convention. The actual leading correction obeys
\begin{equation}
 \boxed{\displaystyle
 D_{M,K}=\mathcal N_0\left\{
 \frac12\log\det G_f-\frac12\log\det G_c
                         -\frac12\log\det Q_\perp\right\}.}
 \label{pqd-main}
\end{equation}
Both scalar determinants are essential. Empty determinants, when there
is no nonbase coarse scalar coordinate, are one.

Here is a direct verification including the measure factors. Let
$K_0$ be the actual retained insertion and $C_p$ its pivot block. The
leading amplitude is
\begin{equation}
 I_{\rm act}=|\det C_p|^{-1}\det A_T^{-1/2},\qquad
 A_T=K_0^T(\mathcal H+T^TT)K_0.\label{pqd-amplitude}
\end{equation}
Its ambient constrained Gaussian is
\begin{equation}
 \int_E dx\,\delta(Cx)
  e^{-\frac12x^T\mathcal H x-\frac12|Tx|^2}
 =(2\pi)^{(\dim E-\dim E_c)/2}I_{\rm act}.
 \label{pqd-gaussian}
\end{equation}
Set $x=U_fu+D_f(\lambda_b,\lambda_c)$. The fine change of variables
has Jacobian $\sqrt{\det G_f}$. The target basis $[U_c,D_c]$ has
Jacobian $\sqrt{\det G_c}$. With
$B=G_c^{-1}D_c^TCU_f$, one has exactly
\begin{equation}
 Cx=U_cAu+D_c(Bu+\lambda_c),\qquad
 \delta(Cx)=(\det G_c)^{-1/2}
                      \delta(Au)\delta(Bu+\lambda_c).
 \label{pqd-target-delta}
\end{equation}
After integrating $\lambda_c=-Bu$, the residual action is
$u^T\mathcal H_\perp u$. Since $TD_b=-I$, integration in
$\lambda_b$ is a translated standard Gaussian and contributes only
$(2\pi)^{\dim S_b/2}$. In particular it produces no background-dependent
pinned-site determinant.

The remaining physical constrained Gaussian, completed by $aA^TA$,
has determinant factor
$a^{\dim\operatorname{ran}P_c/2}\det Q_\perp^{-1/2}$ by
\eqref{pqd-compression}. Since
$\dim S_f=\dim S_b+\dim S_c$, its Gaussian powers agree exactly with
\eqref{pqd-gaussian}. Thus
\begin{equation}
 I_{\rm act}=
 \left(\frac{\det G_f}{\det G_c}\right)^{1/2}
 a^{\dim\operatorname{ran}P_c/2}\det Q_\perp^{-1/2}.
 \label{pqd-amplitude-result}
\end{equation}
The $a$-only power disappears on applying $\mathcal N_0$, proving
\eqref{pqd-main}. The target delta uses a left-relative Lie coordinate.
Its Haar density at the origin is constant; there is no additional
$j(\log W)$ factor.

A useful basis-free ambient form is
\begin{equation}
 Q_{\rm res}=\mathcal H+D_fD_f^T+aC^TP_cC,
 \qquad \det Q_{\rm res}=\det G_f\det Q_\perp,
 \label{pqd-ambient}
\end{equation}
because the full physical and vertical spaces reduce this operator.
On the orthonormal vertical columns $D_fG_f^{-1/2}$ its matrix is
$G_f$. Therefore
\begin{equation}
 D_{M,K}=\mathcal N_0\left\{
 \log\det G_f-\frac12\log\det Q_{\rm res}
                                      -\frac12\log\det G_c\right\}.
 \label{pqd-ambient-identity}
\end{equation}
The physical constraint $P_cC$ is unchanged by $C\mapsto C+D_cZ$.
This observation applies after projection; it does not permit dropping
the endpoint-gauge part of $C$ in the original based ambient precision.

\subsection{A direct canonical Schur check}
\label{pqd-schur-section}

For a second proof, start with the already established canonical
completion
\begin{equation}
 Q_b=\mathcal H+D_bD_b^T+C^TC,\qquad
 D=\mathcal N_0\left(\log\det G_b-\frac12\log\det Q_b\right).
 \label{pqd-canonical}
\end{equation}
Let $J:S_b\to S_f$ be the coordinate insertion, and in orthonormal
full-gauge coordinates set
\begin{equation}
 B_g=G_f^{1/2}J,\qquad C_{gg}=G_c^{1/2}rG_f^{-1/2}.
 \label{pqd-gauge-blocks}
\end{equation}
One has $C_{gg}B_g=0$. Splitting the vertical space into
$\operatorname{ran}B_g$ and its orthogonal complement, the first block
contributes $\det G_b$. In the other vertical block $C_{gg}$ is square
and invertible. Its Schur elimination cancels the whole coarse-gauge
component of $CU_f$, leaving $\mathcal H_\perp+A^TA$. Hence
\begin{equation}
 \det Q_b=\det G_b\det(C_{gg}C_{gg}^T)\det Q_\perp\big|_{a=1}.
 \label{pqd-schur1}
\end{equation}
The scalar Schur identity for \eqref{pqd-scalar-split} gives
$\det(rG_f^{-1}r^T)=\det G_b/\det G_f$. Consequently
\begin{equation}
 \det Q_b=(\det G_b)^2\frac{\det G_c}{\det G_f}
                                      \det Q_\perp\big|_{a=1},
 \label{pqd-schur2}
\end{equation}
which independently checks every determinant power in
\eqref{pqd-main}. Changing $a$ multiplies the physical determinant by
$a^{\dim\operatorname{ran}P_c}$ and changes no normalized expression.

\subsection{The finite identity and the infrared passage are distinct}
\label{pqd-scope}

All preceding identities hold on the actual finite real germ and extend
bilinearly and holomorphically on a sufficiently small complex germ at
that fixed regulator. The ambient form avoids choosing analytic
orthonormal bases. No volume-independent complex neighborhood is
inferred for the globally based scalar inverses.

The actual correction already has a local-density and affine
polynomial-jet interpretation. The separate factors in
\eqref{pqd-ambient-identity} contain massless globally based gauges.
The identity alone does not equate their separate infinite-volume
magnetic traces to the actual Maxwell moment. A joint infrared-regulated
or local-limit identity must retain all three factors. In particular,
one may not replace $G_f,G_c$ separately by massive translation-invariant
ghosts without estimating the change.

\subsection{A uniform flat physical gap}
\label{pqd-flat-section}

This subsection gives estimates for the actual flat operators, uniformly
in $M$, $K=\ell^k$, and compatible volume in $d=4$. The primitive ratio
$\ell$ is fixed. The straight-string angle average is
\begin{equation}
 (A_K^{\rm st}x)_\mu(y)=K^{-d}
 \sum_{z\in B_K(y)}\sum_{j=0}^{K-1}x_\mu(z+je_\mu).
 \label{pqd-angle-average}
\end{equation}
Its absolute row and column sums are $K$ and $K^{1-d}$, respectively,
including periodic multiplicities. Therefore
$\|A_K^{\rm st}\|_{2\to2}\le K^{1-d/2}$.
The actual cumulative ordinary/corrected endpoint-gauge identity gives
$C_K=A_K^{\rm st}+D_c Z_K$ at flatness. Thus
\begin{equation}
 P_cC_K=P_cA_K^{\rm st},\qquad
 \|P_cC_K\|_{2\to2}\le K^{-1}\quad(d=4).
 \label{pqd-projected-bound}
\end{equation}
No bound on $Z_K$ is needed here. The same estimate must not be assigned
to $C_K$ before projection.

For the physical gap, take a flat physical field $x\in\operatorname{ran}P_f$
and put $\eta=K^{-1}$, $v=Kx$ in final physical units. Then
\begin{equation}
 \|v\|_{2,\eta}^2=K^{-2}\|x\|_2^2,\qquad
 P_cC_Kx=P_cQ_K^{\rm st}v.\label{pqd-flat-scaling}
\end{equation}
The four-dimensional action is unchanged by this rescaling. The
previously proved curvature contraction gives, for the actual old-$M$
flat quadratic energy,
\begin{equation}
 \mathcal E_M(v)\ge\|d_\eta v\|_{2,\eta}^2.
 \label{pqd-maxwell-lower}
\end{equation}
The actual coarse Hessian $q_{MK}$ is nonnegative, kills coarse
gradients, and satisfies $\|q_{MK}\|\le R_q$ with a common constant.
The physical residual energy is therefore
\begin{equation}
 \mathcal E_{\rm res}(v)=\mathcal E_M(v)
  -\langle P_cQ_K^{\rm st}v,q_{MK}P_cQ_K^{\rm st}v\rangle\ge0.
 \label{pqd-flat-residual}
\end{equation}

Here is the needed uniform physical Poincare estimate. In each real
coarse Bloch fiber $p$, use fine aliases $p+2\pi r$, gradients
$u_r=K(e^{i(p+2\pi r)/K}-1)$ componentwise, and $P_r=|u_r|^2$.
At flatness globally based and full periodic gradients have the same
range, so $v_r\perp u_r$ is the correct physical condition. Write
$d_c=e^{ip}-1$, and let $\Omega_r$ be the diagonal normalized
straight-average multiplier. The established elementary alias identities
are
\begin{equation}
 P_r\ge4\ (r\ne0),\quad
 \sum_r|\omega_{r,\mu}|^2\le1,\quad
 \min_\mu|\omega_{0,\mu}|\ge c_\omega=(2/\pi)^{d+1},
 \quad \Omega_0u_0=a_0d_c.
 \label{pqd-alias-inputs}
\end{equation}
Set $y=P_{d_c}^{\perp}\sum_r\Omega_rv_r$ and
$y_0=P_{d_c}^{\perp}\Omega_0v_0$. On the zero coarse mode the projection
is $I$, including its harmonic one-forms. For a nonzero mode,
\begin{equation}
 v_0=P_{u_0}^{\perp}\Omega_0^{-1}y_0,
 \qquad |v_0|\le c_\omega^{-1}|y_0|.
 \label{pqd-central-inverse}
\end{equation}
Indeed $\Omega_0v_0-y_0$ is a multiple of $d_c$, whose inverse image
under $\Omega_0$ is a multiple of $u_0$; orthogonal projection removes
it. The zero-mode identity is simply $v_0=y_0$.
Parseval bounds the high-alias contribution to $y$ by $\|v_{\rm high}\|$,
and its energy is at least $4\|v_{\rm high}\|^2$. Consequently
\begin{equation}
 \|v\|^2\le A_0|y|^2+B_0\sum_rP_r|v_r|^2,
 \qquad A_0=2c_\omega^{-2},\quad
 B_0=\frac{1+2c_\omega^{-2}}4.
 \label{pqd-physical-poincare}
\end{equation}
Summing fibers, using \eqref{pqd-maxwell-lower}, and then
\eqref{pqd-flat-residual}, gives
\begin{equation}
 \|v\|_{2,\eta}^2\le B_0\mathcal E_{\rm res}(v)
 +(A_0+B_0R_q)\|P_cQ_K^{\rm st}v\|_2^2.
 \label{pqd-residual-poincare}
\end{equation}
Thus the actual physical precision in \eqref{pqd-physical-precision}
satisfies
\begin{equation}
 \boxed{\displaystyle Q_\perp(1)\ge c_a K^{-2}I,\qquad
 c_a=\min\left\{B_0^{-1},\frac{a}{A_0+B_0R_q}\right\}>0.}
 \label{pqd-flat-gap}
\end{equation}
This is uniform in refinement, endpoint, and finite volume. It supplies
the expected physical gap after the gauge quotient. It does not give
such a gap for the vertical block $G_f$, or bounds on derivatives of
the nonflat projections $P_f,P_c$.

\subsection{A limitation of the unprojected contour completion}
\label{pqd-contour-section}

For clarity, simply returning to the actual contour gauge does not give
the gap \eqref{pqd-flat-gap} for its ambient matrix. First, the full flat
ordinary average has a uniform $O(1)$ counting-$\ell^2$ bound, though
not the $O(K^{-1})$ bound of \eqref{pqd-projected-bound}. Differentiate
the exact cumulative gauge recursion. If $b_\ell$ is the one-step
scalar endpoint contour mean, then, with nested center restrictions,
\begin{equation}
 Z_K=\sum_{j=0}^{k-1}b_\ell^{(j)}A_{\ell^j}^{\rm st}.
 \label{pqd-gauge-recursion}
\end{equation}
A row of $b_\ell$ has absolute norm at most $d\ell$. A row of
$A_r^{\rm st}$ has squared norm at most
$r\,r^{1-d}=r^{2-d}$, by its row sum and largest coefficient. Hence in
$d=4$ each $Z_K$ row has norm at most $d\ell/(1-\ell^{-1})$ and each
$C_K$ row has norm at most
$c_{\rm row}=1+2d\ell/(1-\ell^{-1})$. Actual two-final-block locality
makes each fine bond occur in at most $8d$ final rows. Rowwise
Cauchy--Schwarz therefore proves
\begin{equation}
 \|C_K\|_{2\to2}\le c_C:=\sqrt{8d}\,c_{\rm row}.
 \label{pqd-full-average-bound}
\end{equation}

For $K\ge3$, take one real color and a unit field $e$ on the internal bond $[0,e_1]$
leaving one block center. Put $h=(K-1)/2$ and $\phi=T_Ke$, with zero
values at all block centers and outside this block. In the actual
coordinate-order contour average, for $1\le x_j\le h$ and $j\ge2$,
\begin{equation}
 \phi(1,x_2,\ldots,x_d)=1/d,\qquad
 \phi(0,x_2,\ldots,x_d)=0.\label{pqd-contour-jump}
\end{equation}
The root bond is traversed precisely when direction $1$ is first among
the $d$ nonzero coordinate directions. Set $x=e+D_b\phi$. Exact gauge
identities give $T_Kx=0$, $C_Kx=C_Ke$, and
$x^TL_{\mathrm{flat}}x=e^TL_{\mathrm{flat}}e$. The $h^{d-1}$ bonds in
\eqref{pqd-contour-jump} differ from the original bond and have
$|x|=1/d$, so $\|x\|^2\ge h^{d-1}/d^2$. The common flat operator
bound and \eqref{pqd-full-average-bound} imply
\begin{equation}
 \lambda_{\min}(L_{\mathrm{flat}}+T_K^TT_K+aC_K^TC_K)
 \le d^2(R_q+ac_C^2)h^{-(d-1)}.
 \label{pqd-contour-obstruction}
\end{equation}
In four dimensions this is $O(K^{-3})$, refuting a uniform lower bound
$cK^{-2}$ for that ambient contour completion. It is an upper Rayleigh
bound, not a matching asymptotic. The physical quotient removes exactly
the coordinate-gauge modes used by this test.

\subsection{An infrared-cancelled scalar reference}
\label{pqd-scalar-reference-section}

The following is a specified infinite-lattice reference calculation.
It is not yet a limit theorem for the globally based finite determinants
in \eqref{pqd-main}. Put $\alpha_g=C_H/(48\pi^2)$. The unpinned scalar
magnetic heat coefficient established above is
\begin{equation}
 h(s)=\frac{C_H}{3}s^2e^{-8s}I_0(2s)^2I_1(2s)^2,
 \quad h(s)=O(s^4)\ (s\downarrow0),\quad
 h(s)=\alpha_g+O(s^{-1})\ (s\to\infty).
 \label{pqd-scalar-heat}
\end{equation}
In common final physical heat time $u$, the fine and coarse coefficients
are $h(K^2u)$ and $h(u)$. Their orbit-metric reference ratio is exactly
\begin{equation}
 \frac12\int_0^\infty\{h(K^2u)-h(u)\}\frac{du}{u}
                       =\alpha_g\log K.
 \label{pqd-scalar-scale}
\end{equation}
The integral is absolutely convergent. On $\epsilon<u<T$ its unhalved
value equals
$\int_T^{K^2T}h(v)\,dv/v-\int_\epsilon^{K^2\epsilon}h(v)\,dv/v$;
the endpoint limits prove \eqref{pqd-scalar-scale}. A common physical
mass $\delta$ multiplies the integrand by $e^{-\delta^2u}$, and dominated
convergence gives the same limit as $\delta\downarrow0$, for each
fixed $K$.

This exact scalar cancellation is useful comparison data. Applying it
to the actual correction still requires a joint infrared passage for
\eqref{pqd-main}, followed by uniform curvature-source estimates for the
projected residual operator. The flat physical gap alone supplies
neither assertion. Conditional on those further bridges, the physical
term $-\frac12\log\det Q_\perp$ would have to contribute
$21\alpha_g\log K$ to obtain the target total
$22\alpha_g\log K$; that contribution is not proved here.



\section{The remaining coupled estimate}

Let
\[
 B_v(M,K)=-[t^2]\operatorname{Tr}_{0,E}\log V_a(t).
\]
The exact determinant identity, the bounded coarse
contribution and Theorem~\ref{pgl-full} give
\[
 \beta(D_{M,K})=B_v(M,K)+4\alpha_g\log K+O(1).
\]
The specified massive continuum reference has total
coefficient \(22\alpha_g\log K+O(1)\).
Therefore a uniform bounded comparison with that
reference is equivalent to the remaining estimate
\[
                  B_v(M,K)=18\alpha_g\log K+O(1)
\]
over the full required family of preceding depths and
structural systems.

If the actual vector prefix through time \(K^2\)
has its continuum value \(20\alpha_g\log K+O(1)\),
its remaining tail must be
\(-2\alpha_g\log K+O(1)\).
This last statement is conditional on that vector
prefix estimate. The scalar calculation alone gives
no sign for the actual combined determinant.
In particular a proof based on separately bounded
vector and ghost tails past \(K^2\) cannot work:
the ghost tail has a proved logarithmic lower bound.

The full actual short-time interval \(0<s<1\) is
now bounded by Theorem~\ref{vst-theorem} and the
ghost estimates above. The unresolved analytic work
lies at intermediate and longer times; it must
retain the cancellation just identified.

The finite physical-quotient identity
\eqref{pqd-main} provides another route to that
combined estimate. It eliminates periodic pinning
algebraically, keeps the fine/coarse orbit-density
ratio, and exposes the flat physical gap
\eqref{pqd-flat-gap}. The next work on this route is
to justify the joint infrared passage with all
factors retained and to control the physical
precision's curvature derivatives. Its spectral
gap alone does not establish the required locality
or the determinant coefficient.

Once a positive actual block margin is established
with the needed uniformity, the finite drift-transfer
and finite shooting results of Notebook 32 apply.
The global measure matching, mixed and large-field
histories, sourced continuum limits, nontriviality,
Euclidean covariance, positivity and reconstruction
remain requirements of the full existence objective.

\begin{thebibliography}{9}
\bibitem{priorpg}
\emph{Toward a four-dimensional Yang--Mills construction},
Notebooks 32 and 33 in this workspace.
The actual determinant normalization, individually
gauge-invariant densities, affine stationary jets
and proper-time reference are used with their
stated scope.
\bibitem{sourcepg}
E.~Sussman and collaborators,
\href{https://github.com/ethanwsussman/Yang--Mills-LLM-experiment}
{Yang--Mills-LLM-experiment}, deposited accounts pinned at
commit \texttt{e392ef45013a581664bab76b90867170b30fe10e}.
\end{thebibliography}
\end{document}
